\section{Reproducibility Contract}\label{app:reproducibility}

The Python package is located under \texttt{src/}.  From a clean checkout, the
intended local setup is:

\begin{verbatim}
python -m venv .venv
.venv/Scripts/python -m pip install -e "src[dev]"
.venv/Scripts/python -m pytest -q -p no:cacheprovider src/a2a_detection/tests
\end{verbatim}

On POSIX systems the interpreter path is \texttt{.venv/bin/python}.  No PyPI
package or public repository URL is asserted by this manuscript.  A test pass
verifies software contracts only; it does not validate the scientific
construct or the missing data.

The audit entry point, run from the repository root, is:

\begin{verbatim}
PYTHONPATH=src python -B -m a2a_detection.scripts.audit_research_claims \
  --strict --output GPD/verification/research-audit-results.json
\end{verbatim}

Input paths default to the historical files under \texttt{data/} and can be
overridden with \texttt{--scores}, \texttt{--labels},
\texttt{--raw-transactions}, and \texttt{--injection}.  With
\texttt{--strict} the command writes the JSON and then exits with status 1
whenever any section is invalid, stale, or missing; on the current historical
artifacts it therefore exits 1 by design.  Without \texttt{--strict} it writes
the same JSON and exits 0.  It must fail closed when required artifacts, event
indexes, schemas, or hashes are absent.  The historical ignored files are not sufficient for a clean-clone
reproduction, so the empirical reproducibility gate remains failed even if the
local script can reconstruct their statistics.

Every future release must add one machine-readable manifest containing the
dataset version, SHA-256 hashes, query cutoff, package revision, Python and
dependency versions, random seeds, split identifiers, model configuration,
threshold-selection rule, and generated table values.  Paper tables should be
rendered from that output or checked against it in continuous integration.
